\documentclass[a4paper,12pt]{article}
\setcounter{secnumdepth}{5}
\setcounter{tocdepth}{3}
\input{/usr/share/LaTeX-ToolKit/template.tex}
\renewcommand{\arraystretch}{1.5}
\begin{document}
\title{計算機程式期末專題期末報告}
\author{沈威宇}
\date{\temtoday}
\titletocdoc
\sct{第一部分：專題整體}
\sct{第二部分：個人程式設計說明}
\ssc{總論}
\sssc{我撰寫的部分}
\bit
\item \verb|include/algorithm-copy.hpp|：全部
\item \verb|include/array.hpp|：全部
\item \verb|include/range-access.hpp|：全部
\item \verb|include/algorithm-find.hpp|：全部
\item \verb|include/circulate.hpp|：全部
\item \verb|include/swap.hpp|：全部
\item \verb|include/algorithm-heap.hpp|：全部
\item \verb|include/ignore.hpp|：全部
\item \verb|include/tuple.hpp|：全部
\item \verb|include/algorithm-remove.hpp|：全部
\item \verb|include/priority_queue.hpp|：全部
\item \verb|include/tuple-using.hpp|：全部
\item \verb|include/allocator.hpp|：全部
\item \verb|include/qsort.hpp|：全部
\item \verb|include/vector.hpp|：全部
\item \verb|Mods.hpp|：全部
\item \verb|mods/GameOfLife.hpp|：全部
\item \verb|Game.hpp|：全部
\item \verb|Renderer.hpp|：全部
\item \verb|generate-notes.hpp|：全部
\item \verb|KeyNoteData.hpp|：全部
\item \verb|RhythmQuest.cpp|：大部分
\eit
\sssc{報告架構}
在本章節後續內容中，我將依照課程要求之 A、B、C、D 四項題目（如適用），分別以模組或檔案群組為單位進行說明。
\ssc{include/algorithm-copy.hpp}
The \verb|algorithm-copy.hpp| file provides three basic data transfer algorithms: \verb|copy|, \verb|copy_if|, and \verb|copy_backward|. This file serves as the algorithm module for this project's custom standard function library, mimicking the interface design of C++20 \verb|<algorithm>|. The copy function family demonstrates the polymorphism programming techniques, particularly conceptual parameter constraints and conditional optimization strategies.

I utilize concepts for constraints, such as \verb|std::input_iterator|, \verb|std::weakly_incrementable|, \verb|std::bidirectional_iterator|, and \verb|std::indirectly_writable| in the definition, and
\verb|std::contiguous_iterator| and \verb|std::is_trivially_copyable_v| for specialized acceleration.

In \verb|copy| algorithms, I use \verb|std::memcpy| when possible and element-wise assignment only when necessary, as demonstrated below (excerpted from \verb|include/algorithm-copy.hpp| lines 7-32)
\begin{lstlisting}[language=C++]
template <std::input_iterator InputIt, std::weakly_incrementable OutputIt>
  requires std::indirectly_writable<OutputIt, std::iter_value_t<InputIt>>
constexpr OutputIt copy(InputIt first, InputIt last, OutputIt d_first) {
  if constexpr (std::contiguous_iterator<InputIt> &&
                std::contiguous_iterator<OutputIt> &&
                std::is_trivially_copyable_v<
                    typename std::iterator_traits<InputIt>::value_type> &&
                std::is_same_v<
                    typename std::iterator_traits<InputIt>::value_type,
                    typename std::iterator_traits<OutputIt>::value_type>) {
    if (first == last)
      return d_first;
    const auto *src = std::to_address(first);
    const auto *src_end = std::to_address(last);
    auto *dest = std::to_address(d_first);
    const std::size_t count = (src_end - src) * sizeof(*src);
    std::memcpy(dest, src, count);
    return d_first + (last - first);
  }
  while (first != last) {
    *d_first = *first;
    ++first;
    ++d_first;
  }
  return d_first;
}
\end{lstlisting}

Similarly, in \verb|copy_backward| algorithm, \verb|std::memmove| is used when possible, as demonstrated below (excerpted from \verb|include/algorithm-copy.hpp| lines 49-73)
\begin{lstlisting}[language=C++]
template <std::bidirectional_iterator BidirIt1,
          std::bidirectional_iterator BidirIt2>
constexpr BidirIt2 copy_backward(BidirIt1 first, BidirIt1 last,
                                 BidirIt2 d_last) {
  using ValueType1 = typename std::iterator_traits<BidirIt1>::value_type;
  using ValueType2 = typename std::iterator_traits<BidirIt2>::value_type;
  if constexpr (std::contiguous_iterator<BidirIt1> &&
                std::contiguous_iterator<BidirIt2> &&
                std::is_trivially_copyable_v<ValueType1> &&
                std::is_same_v<ValueType1, ValueType2>) {
    auto n = last - first;
    if (n > 0) {
      std::memmove(std::addressof(*d_last) - n, std::addressof(*first),
                   n * sizeof(ValueType1));
    }
    return d_last - n;
  } else {
    while (first != last) {
      --last;
      --d_last;
      *d_last = *last;
    }
    return d_last;
  }
}
\end{lstlisting}
\subsection{include/allocator.hpp}

The file \texttt{allocator.hpp} implements a custom memory management system for the project's containers, including \texttt{vector}, \texttt{priority\_queue}, and other generic data structures. It defines both the \texttt{allocator} class and the \texttt{allocator\_traits} class, along with helper templates and utilities to enable allocator-aware construction and deallocation, following C++20 standard.
\subsubsection*{題目 A — Class and Module Structure}
The allocator module consists of the following main components:
\begin{itemize}
  \item \textbf{\texttt{mystd::allocator<T>}}: A template class that implements memory allocation, deallocation, and default construction/destruction for type \texttt{T}. Key features include:
    \begin{itemize}
      \item \texttt{allocate(n)}: Allocates memory for \texttt{n} objects of type \texttt{T}, with overflow checks.
      \item \texttt{deallocate(p, n)}: Deallocates memory previously allocated for \texttt{n} objects at pointer \texttt{p}.
      \item Default constructors, copy constructors, and destructor are trivial or \texttt{noexcept}.
      \item Supports equality comparison between different allocator types.
      \item Provides type aliases like \texttt{value\_type}, \texttt{size\_type}, \texttt{difference\_type}, and propagation traits.
    \end{itemize}
  \item \textbf{\texttt{mystd::allocator\_traits<Alloc>}}: A template struct that encapsulates allocator traits and provides generic interfaces to interact with any allocator type. Key responsibilities:
    \begin{itemize}
      \item Deduce \texttt{pointer}, \texttt{const\_pointer}, \texttt{void\_pointer}, \texttt{const\_void\_pointer}, \texttt{size\_type}, and \texttt{difference\_type} from \texttt{Alloc}.
      \item Provide facilities for \texttt{allocate} and \texttt{deallocate}, optionally with a memory hint.
      \item Provide construction and destruction methods that respect allocator capabilities: \texttt{construct(a, p, args...)} and \texttt{destroy(a, p)}.
      \item Provide container propagation flags: \texttt{propagate\_on\_container\_copy\_assignment}, \texttt{propagate\_on\_container\_move\_assignment}, \texttt{propagate\_on\_container\_swap}, and \texttt{is\_always\_equal}.
      \item Support allocator rebinds: \texttt{rebind\_alloc<U>} and \texttt{rebind\_traits<U>}.
    \end{itemize}
  \item \textbf{Helper Templates}:
    \begin{itemize}
      \item \texttt{\_\_rebind\_helper<Alloc, T>}: Determines the allocator type rebound to a different element type \texttt{T}, supporting both classic \texttt{rebind} and variadic template allocators.
      \item Pointer type helpers (\texttt{pointer\_helper}, \texttt{const\_pointer\_helper}, etc.): Deduce correct pointer types, with SFINAE detection for user-defined types.
      \item Compile-time detection helpers (\texttt{has\_construct}, \texttt{has\_destroy}, \texttt{test}): Detects whether an allocator provides specialized construct/destroy interfaces.
    \end{itemize}
  \item \textbf{Allocator-aware Object Construction Utilities}:
    \begin{itemize}
      \item \texttt{uses\_allocator\_construction\_args}: Computes the proper argument tuple for constructing an object of type \texttt{T} with a given allocator.
      \item \texttt{make\_obj\_using\_allocator}: Constructs a new object using the allocator-aware arguments tuple.
      \item \texttt{uninitialized\_construct\_using\_allocator}: Constructs an object at a pre-allocated memory location using allocator-aware construction.
      \item Traits and utility functions for \texttt{uses\_allocator} and allocator-aware construction, supporting tuples, piecewise construction, and pairs.
    \end{itemize}
\end{itemize}
\subsubsection*{題目 B — Constructors, Destructors, and Object Management}

\begin{itemize}
  \item The \texttt{allocator} class provides default constructors and destructors that are trivial and \texttt{noexcept}.
  \item Allocation is performed via \texttt{::operator new}, with overflow protection for large object counts.
  \item Deallocation is performed via \texttt{::operator delete}, ensuring memory safety.
  \item Object construction uses either the allocator’s \texttt{construct} method (if available) or \texttt{std::construct\_at}, determined at compile-time via SFINAE and \texttt{if constexpr}.
  \item Object destruction similarly uses either \texttt{destroy} from the allocator or \texttt{std::destroy\_at}.
  \item Rebinding support allows the allocator to be reused for different element types in generic containers.
\end{itemize}

\subsubsection*{題目 C — Object-Oriented and Modern C++ Techniques}

\begin{itemize}
  \item \textbf{Template Meta-programming}: Uses SFINAE (\texttt{std::void\_t}) and compile-time detection to adapt behavior depending on allocator capabilities.
  \item \textbf{Concept-based Type Checking}: Ensures generic allocator functions work with any type satisfying allocator requirements.
  \item \textbf{Tuple-based Forwarding for Allocator-aware Construction}: Correctly forwards arguments to constructors depending on whether the type supports \texttt{uses\_allocator}.
  \item \textbf{Piecewise Construction Support}: For pair types, supports tuple-based piecewise construction.
  \item \textbf{Compile-time Optimizations}: Uses \texttt{if constexpr} to select the appropriate construction/destruction path without runtime overhead.
\end{itemize}

\begin{verbatim}
// Example of allocator-aware construction (simplified)
auto args = mystd::uses_allocator_construction_args<MyType>(alloc, arg1, arg2);
std::apply([&](auto&&... xs){ std::construct_at(p, std::forward<decltype(xs)>(xs)...); }, args);
\end{verbatim}

\subsubsection*{題目 D — Special Techniques and Performance Optimizations}

\begin{itemize}
  \item Zero-overhead abstraction for allocator traits and memory management.
  \item Compile-time detection of allocator capabilities prevents unnecessary runtime checks.
  \item Propagation flags reduce unnecessary copies or allocations when containers are copied, moved, or swapped.
  \item Efficient construction and destruction mechanisms via tuple forwarding and SFINAE minimize unnecessary object copies.
  \item Integrates seamlessly with project containers to maintain memory safety, correct alignment, and high performance for large-scale container operations.
  \item Supports generic programming and software reuse across multiple container types without sacrificing efficiency.
\end{itemize}


